\documentclass[10pt,a4paper]{amsart}
\usepackage[margin=19mm]{geometry}
\usepackage{amsmath,amssymb,mathtools}
\usepackage[scr=rsfs,cal=euler]{mathalfa}
\usepackage{xfrac}
\usepackage[unicode,hidelinks,bookmarks=false]{hyperref}
\newcommand{\ie}{i.e.\ }
\newcommand{\cf}{cf.\ }
\newcommand{\resp}{respectively\ }
\newcommand{\Subsection}{\subsection}
\providecommand{\nobreakdash}{\nobreak\hbox{-}}
\setlength{\emergencystretch}{2em}
\multlinegap0pt
\usepackage{lmodern}
\usepackage{mathtools}
\usepackage{needspace}
\hypersetup{pdftitle={On Kolmogorov's rearrangement problem and Garsia's conjecture},pdfauthor={Mark Lewko}}
\newtheorem{theorem}{Theorem}
\newtheorem{proposition}[theorem]{Proposition}
\newtheorem{lemma}[theorem]{Lemma}
\newtheorem{corollary}[theorem]{Corollary}
\newtheorem{remark}[theorem]{Remark}
\numberwithin{equation}{section}
\newcommand{\T}{\mathbb T}
\newcommand{\N}{\mathbb N}
\newcommand{\Z}{\mathbb Z}
\newcommand{\R}{\mathbb R}
\newcommand{\C}{\mathbb C}
\newcommand{\E}{\mathbb E}
\newcommand{\PP}{\mathbb P}
\newcommand{\one}{\mathbf 1}
\newcommand{\calM}{\mathcal M}
\newcommand{\calF}{\mathcal F}
\newcommand{\calA}{\mathcal A}
\newcommand{\calB}{\mathcal B}
\newcommand{\calE}{\mathcal E}
\newcommand{\norm}[1]{\left\lVert#1\right\rVert}
\setlength{\emergencystretch}{2em}
\allowdisplaybreaks[2]
\begin{document}
\title[On Kolmogorov's rearrangement problem and Garsia's conjectur]{On Kolmogorov's rearrangement problem and Garsia's conjecture}
\author{Mark Lewko}
\date{}
\maketitle
\noindent\emph{Source and licence.} On Kolmogorov's rearrangement problem and Garsia's conjecture. Version arXiv:2609.18491v1 (CC BY 4.0). This material is distributed under the Creative Commons Attribution 4.0 International licence, http://creativecommons.org/licenses/by/4.0/, as recorded in the arXiv OAI licence metadata for this exact version and shown on the abstract page https://arxiv.org/abs/2609.18491v1. Full rights notice: licenses/arxiv-2609.18491-CC-BY-4.0.txt.\par\smallskip
\noindent\textit{Editorial scope.} Counted unit: the original \texttt{lemma} environment \texttt{lem:two} at source lines 99--110 together with its complete proof at lines 170--196; the delivered document keeps source lines 87--197 so that every referenced label and every citation resolves inside it. Native editable source is the paper's own concatenated TeX; the AMS wrapper is editorial formatting only. No publisher class, upstream script or unrelated artwork is redistributed.
\begin{equation}\label{eq:construction}
 \phi_n(x,0):=e(nx),\qquad
 \phi_n(x,1):=e(\sigma(n)x),\qquad n\in[N],
\end{equation}
on $\T\times\{0,1\}$, where each copy of $\T$ has measure $1/2$. These functions are orthonormal and have absolute value one. We choose $\sigma$ using the combinatorial lemma in the next section, then apply the classical divergence theorem for rearranged Fourier series.

\section{The combinatorial argument}

Translating the frequencies of a trigonometric polynomial by an integer, or dilating them by a positive integer, preserves the distribution of its maximal absolute partial sum. A finite Fourier example can therefore be placed on any arithmetic progression, provided that the permutation of its frequencies is preserved. This is why arithmetic progressions enter the construction.

Fix a positive integer $m$ and a permutation $\pi$ of $[m]$. In a permuted list of integers, we seek an arithmetic progression $x_1<\cdots<x_m$ whose terms occur as $x_{\pi(1)},\ldots,x_{\pi(m)}$. Other integers in the list may occur between these terms. The following lemma provides the permutation needed for the two-copy construction.


\begin{lemma}\label{lem:two}
Let $\pi$ be a permutation of $[m]$. There exist $N\geq m$ and a permutation $\sigma$ of $[N]$ such that, for every permutation $\tau$ of $[N]$, there is an $m$-term arithmetic progression $x_1<\cdots<x_m$ in $[N]$ for which
\[
 x_{\pi(1)},\ldots,x_{\pi(m)}
\]
is a subsequence of at least one of
\[
 \tau(1),\ldots,\tau(N)
 \qquad\text{and}\qquad
 \sigma(\tau(1)),\ldots,\sigma(\tau(N)).
\]
\end{lemma}


For the system \eqref{eq:construction}, a permutation $\tau$ of the functions gives the frequency permutations $\tau$ and $\sigma\circ\tau$ on the two copies of $\T$.

We prove the lemma in two steps: first count the colorings that fail to contain the required progression, then use this count to choose $\sigma$. Here $m$ remains the progression length. We introduce $K\geq m$ available colors, labeled $1,\ldots,K$; the chosen progression may use any $m$ distinct colors. The counting estimate holds for every fixed $K\geq m$, and the final step will use $K=m^2$.

\begin{lemma}\label{lem:count}
Fix integers $m\geq2$ and $K\geq m$, and a permutation $\pi$ of $[m]$. All but at most
\begin{equation}\label{eq:count}
 (m-1)^N\exp(o(N))
\end{equation}
of the $K^N$ colorings $c:[N]\to[K]$ admit an $m$-term arithmetic progression $x_1<\cdots<x_m$ in $[N]$ with distinct colors such that $x_{\pi(j)}$ has the $j$-th smallest color among these $m$ points, for each $j\in[m]$. Equivalently,
\begin{equation}\label{eq:pattern}
 c(x_{\pi(1)})<\cdots<c(x_{\pi(m)}).
\end{equation}
Here $N\to\infty$ with $m$, $K$, and $\pi$ fixed.
\end{lemma}

For the application to Lemma~\ref{lem:two}, where $K=m^2$, the point is that the square of this bound is exponentially smaller than $K^N=m^{2N}$. The bound is sharp up to the $\exp(o(N))$ factor, since a coloring using at most $m-1$ colors admits no progression with $m$ distinct colors.

\begin{proof}
Let $\calF_N$ be the family of colorings $c:[N]\to[K]$ containing no progression satisfying \eqref{eq:pattern}. We count these colorings by deleting the points $N,N-1,\ldots,1$ in turn.

\emph{The counting step.}
Consider any family $\calA$ of colorings of $[r]$, where $1\leq r\leq N$. We delete the last point $r$. For each color $i\in[K]$, let
\[
 \calA_i:=\{c|_{[r-1]}:c\in\calA,\ c(r)=i\}.
\]
Here $[0]:=\varnothing$. Thus $\calA_i$ consists of the colorings of $[r-1]$ that extend to a member of $\calA$ by giving $r$ color $i$. We claim that
\begin{equation}\label{eq:recursion}
 |\calA|\leq(m-1)\left|\bigcup_{i=1}^K\calA_i\right|
 +\sum_{\substack{Q\subseteq[K]\\|Q|=m}}
 \left|\bigcap_{i\in Q}\calA_i\right|.
\end{equation}
Indeed, a coloring of $[r-1]$ that belongs to exactly $s$ of the families $\calA_i$ has $s$ extensions in $\calA$. If $s\geq1$, it contributes $m-1+\binom sm$ on the right, and this is at least $s$: for $s\leq m$ this is clear, and for $s>m$ we have $\binom sm\geq\binom s1=s$. If $s=0$, it contributes zero to both sides.

Apply \eqref{eq:recursion} first with $\calA=\calF_N$ and $r=N$, then to each family on the right with $r=N-1$, and continue. At each application, $\calA$ denotes the family currently being counted, and $\calA_i$ is defined from that family by deleting its last point. A union contributes a factor $m-1$; an intersection contributes a choice of an $m$-element color set $Q$.

\Needspace{6\baselineskip}
\emph{Controlling the intersections.}
Fix one term in the resulting sum. Let $S\subseteq[N]$ be the points at which an intersection was chosen, and let $Q_x$ be the color set chosen at $x\in S$. These choices determine the term. Once all points have been deleted, the remaining family is either empty or consists of the unique coloring of the empty set. A nonzero term therefore has weight $(m-1)^{N-|S|}$.

For such a term, we may prescribe any color from $Q_x$ at each $x\in S$ and extend these choices to a coloring in $\calF_N$. To see this, restore the points in reverse order. Membership in a union allows at least one color at the restored point; membership in the chosen intersection allows every color in $Q_x$, so we use the prescribed one. The colors outside $S$ may depend on these prescriptions.

Fix an $m$-element color set $Q=\{q_1<\cdots<q_m\}$. The points $x\in S$ with $Q_x=Q$ contain no $m$-term arithmetic progression. Otherwise, on such a progression $x_1<\cdots<x_m$, prescribe color $q_j$ at $x_{\pi(j)}$. The extension just described would belong to $\calF_N$ and satisfy \eqref{eq:pattern}, a contradiction.

Let $r_m(N)$ denote the largest size of a subset of $[N]$ containing no $m$-term arithmetic progression. By Szemer\'edi's theorem \cite{Szemeredi}, $r_m(N)=o(N)$. Each of the $\binom Km$ classes of points with the same color set has size at most $r_m(N)$, and hence
\[
 |S|\leq\binom Km r_m(N)=o(N)
\]
uniformly over the nonzero terms. Let $d_N$ be the largest possible $|S|$ among these terms; thus $d_N=o(N)$.

For $|S|=j$, there are $\binom Nj$ choices of $S$ and $\binom{K}{m}^{j}$ choices of its color sets. Summing the weights gives
\[
 |\calF_N|\leq\sum_{j=0}^{d_N}\binom Nj\binom{K}{m}^{j}(m-1)^{N-j}
 \leq(m-1)^N\exp(o(N)),
\]
where the last estimate uses $d_N=o(N)$.
\end{proof}


\begin{proof}[Proof of Lemma~\ref{lem:two}]
For $m=1$, take $N=1$. For $m\geq2$, set $K:=m^2$ and take $N$ divisible by $K$. Call a coloring $c:[N]\to[K]$ \emph{balanced} if each color is used exactly $N/K$ times, and let $\calB_N$ be the family of balanced colorings. Then
\[
 |\calB_N|=:T_N=\frac{N!}{((N/K)!)^K}=K^N\exp(o(N)).
\]
Let $\calF_N$ be the exceptional family in Lemma~\ref{lem:count}. For a fixed balanced coloring $c$ and a uniformly random permutation $\sigma$ of $[N]$, the coloring $c\circ\sigma^{-1}$ is uniform on $\calB_N$. Thus the expected number of balanced colorings $c$ for which both $c$ and $c\circ\sigma^{-1}$ belong to $\calF_N$ is
\[
 \frac{|\calF_N\cap\calB_N|^2}{T_N}
 \leq\frac{|\calF_N|^2}{T_N}
 \leq\left(\frac{(m-1)^2}{m^2}\right)^N\exp(o(N))\longrightarrow0.
\]
For $N$ sufficiently large, choose $\sigma$ for which there is no such coloring.

Now fix an arbitrary permutation $\tau$ of $[N]$. Partition the positions in the sequence $\tau(1),\ldots,\tau(N)$ into $K$ consecutive intervals, each of length $N/K$. Define the coloring $c_\tau:[N]\to[K]$ by
\begin{equation}\label{eq:block-coloring}
 c_\tau(\tau(j))=\ell
 \quad\text{when}\quad
 (\ell-1)\frac{N}{K}<j\leq\ell\frac{N}{K},
 \qquad 1\leq\ell\leq K.
\end{equation}
Thus the color of an integer records the interval in which it appears in this particular permutation. Define $c_{\sigma\circ\tau}$ by the same rule for the sequence $\sigma(\tau(1)),\ldots,\sigma(\tau(N))$. These colorings are balanced, and
\[
 c_{\sigma\circ\tau}(\sigma(n))=c_\tau(n),
 \qquad n\in[N].
\]
In particular, $c_{\sigma\circ\tau}=c_\tau\circ\sigma^{-1}$. By the choice of $\sigma$, at least one of $c_\tau$ and $c_{\sigma\circ\tau}$ lies outside $\calF_N$. Lemma~\ref{lem:count} supplies a progression whose terms, permuted by $\pi$, have strictly increasing colors for that coloring. By \eqref{eq:block-coloring}, smaller colors correspond to earlier positions in the corresponding permutation. Hence $x_{\pi(1)},\ldots,x_{\pi(m)}$ is the required subsequence.
\end{proof}

\begin{thebibliography}{99}
\bibitem[Szemeredi]{Szemeredi} E. Szemer\'edi, \emph{On sets of integers containing no $k$ elements in arithmetic progression}, Acta Arith. \textbf{27} (1975), 199--245. \href{https://doi.org/10.4064/aa-27-1-199-245}{doi:10.4064/aa-27-1-199-245}.
\end{thebibliography}
\end{document}
